\section{Trading degree for fewer direct proof coordinates}
\label{sec:monic-elimination}

The preceding sign bootstrap gives a direct equation in $248$ unrestricted
integers.  We now reduce its dimension by allowing higher degree.  This
route retains the direct C222 proof calculus: no K215 simulation, Jones
universal equation, or large fixed coefficient tuple is introduced.
The elimination principle is standard.  The contribution here is the
complete family of explicit polynomials and the certified tradeoffs below.

Write the preceding source as
\[
 P_{248}(x)=\sum_{j=0}^{181}R_j(x)^2,
 \qquad x=(x_0,\ldots,x_{247})\in\Z^{248}.
\]
The complete quadratic residual list is retained without change in
\texttt{source/equation248.json}; its source hash is recorded in
\texttt{SOURCE.json}.  The pivot tables in
Appendix~\ref{app:monic-pivots}, together with that literal residual list,
specify every new polynomial by finite substitution.

\begin{theorem}[Acyclic monic elimination]
\label{thm:acyclic-monic}
Let $R_0,\ldots,R_{m-1}\in\Z[x_0,\ldots,x_{N-1}]$.
Select different rows and different coordinates so that each selected row
has the form
\[
 R_{r(v)}=\varepsilon_v x_v+f_v,
 \qquad \varepsilon_v\in\{-1,1\},
 \qquad x_v\text{ does not occur in }f_v.
\]
Suppose the directed graph on selected coordinates, with an edge $v\to w$
when $x_w$ occurs in $f_v$, is acyclic.  List the retained source indices
as $i_0<\cdots<i_{n-1}$, and let $y=(y_0,\ldots,y_{n-1})$.
Define integer polynomials $F_v(y)$ by $F_{i_j}=y_j$ on retained
coordinates, and recursively by
\[
 F_v=-\varepsilon_v f_v(F)
\]
on selected coordinates.  If $I$ is the set of unselected rows, put
\[
 Q(y)=\sum_{j\in I}R_j(F(y))^2.
\]
Projection to retained coordinates is a bijection between the integer
zero sets of $\sum_jR_j^2$ and $Q$.  Its inverse is the polynomial map
$y\mapsto(F_0(y),\ldots,F_{N-1}(y))$.
\end{theorem}

\begin{proof}
At an integer zero of the source sum of squares, every residual is zero.
Each selected row therefore forces exactly
$x_v=-\varepsilon_v f_v$.  Induction along a reverse topological order of the
acyclic graph gives $x_v=F_v(y)$ for every coordinate.  The unselected
residuals vanish, so $Q(y)=0$.

Conversely, if $Q(y)=0$, every remaining squared integer residual is zero.
Set $x=F(y)$.  The selected residuals vanish by their definitions, and
the unselected ones by $Q(y)=0$.  Thus the source polynomial is zero.
Every restored coordinate is an integer because all the substitutions
have integer coefficients.  The same induction proves uniqueness.
No division, nonzero denominator, congruence, or sign choice is needed.
\end{proof}

The unrestricted integer domain matters.  Applying this theorem to a
nonnegative-domain equation would require additional proofs that all
restored coordinates remain nonnegative at every zero.  We use the
integer source throughout.

\subsection{The resulting explicit family}

Every row in the following table defines one ordinary polynomial.  Its
complete arithmetic DAG uses only integer constants, its stated input
variables, and binary addition, subtraction and multiplication.
Intermediate registers abbreviate polynomials and are not additional
quantified unknowns.  Degree is exact.  The last column is a proved bound
on the height after collecting all coefficients; it is not the largest
constant appearing in the DAG.

\begin{center}
\begin{tabular}{rrrr}
\hline
Integer unknowns & Remaining residuals & Degree & Height bound \\
\hline
225 & 159 & 6 & $<2^{27}$ \\
153 & 87 & 8 & $<2^{47}$ \\
144 & 78 & 12 & $<2^{50}$ \\
125 & 59 & 16 & $<2^{80}$ \\
117 & 51 & 32 & $<2^{144}$ \\
111 & 45 & 64 & $<2^{173}$ \\
102 & 36 & 256 & $<2^{691}$ \\
94 & 28 & $65{,}540$ & $<2^{164870}$ \\
\hline
\end{tabular}
\end{center}

The degree-$16$ equation removes $123$ of the source's $248$ coordinates,
or $49.6\%$.  The degree-$256$ equation removes $146$, or $58.9\%$.
These are alternatives within the direct proof-calculus route, not
global dimension records among all coefficient encodings.  The
$94$-variable equation illustrates the cost of prioritizing dimension
alone: its degree and coefficient bound are much larger.

For the $153$-variable octic the package also supplies every collected
coefficient.  Its exact expansion has $188{,}579$ nonzero monomials and
height
\[
 750{,}273{,}658{,}864.
\]
Its $87$ fully substituted residuals contain $3{,}476$ monomials in total.
This exact height is stronger than the general $2^{47}$ bound.  The
expansion makes the support cost of eliminating coordinates explicit.

\subsection{Exact degree and coefficient bounds}

Give each retained coordinate degree one.  For each restored coordinate
propagate the maximum, over its monomials, of the sum of the previously
assigned coordinate degrees.  Apply the same rule to each remaining
residual.  If the largest bound is $d$, then $\deg Q\le2d$.

To prove equality, the supplied certificate gives a remaining residual
of degree bound $d$ whose degree-$d$ homogeneous part is nonzero when
the retained coordinate with source index $i$ is set to $i+3$, reduced
modulo $1{,}000{,}000{,}007$.  A nonzero modular value proves a nonzero
integer value; no primality assumption is needed.  At this integer point
the sum of the squares of all degree-$d$ homogeneous parts is positive.
It follows that $\deg Q=2d$.  This is an exact nonvanishing certificate,
not probabilistic polynomial identity testing.

\begin{lemma}[A propagated height bound]
\label{lem:monic-height}
Set $L_v=1$ for each retained coordinate.  For a restored polynomial
$F_v=\sum_M c_M\prod_{j\in M}F_j$, set
\[
 L_v=\sum_M |c_M|\prod_{j\in M}L_j.
\]
Use the same rule to obtain a bound $L_j'$ for each remaining residual.
Then every coefficient of $Q$ has absolute value at most
\[
 B_Q=\sum_{j\in I}(L_j')^2.
\]
If this positive integer has bit length $b$, every coefficient has
absolute value strictly below $2^b$.
\end{lemma}

\begin{proof}
For an integer polynomial $f$, let $\|f\|_1$ be the sum of the absolute
values of its collected coefficients.  The triangle and convolution
inequalities give
\[
 \|f+g\|_1\le\|f\|_1+\|g\|_1,
 \qquad \|fg\|_1\le\|f\|_1\|g\|_1.
\]
Induction through the substitutions proves $\|F_v\|_1\le L_v$.
Applying the same inequalities to the residuals and their squared sum
proves $\|Q\|_1\le B_Q$.  Each coefficient is bounded by this norm,
and $B_Q<2^b$ by the definition of bit length.
\end{proof}

All stated bit bounds are obtained by exact integer arithmetic.  The
largest literal constant in these DAGs is $256$; that different statistic
does not replace the expanded coefficient bounds in the table.

\subsection{Verification and the consistency theorem}

The independent checker expands each circuit block as a polynomial in
the coordinates already available.  These local expansions have
degree at most two, even when the fully substituted polynomial has degree
$65{,}540$.  It checks the blocks coefficient by coefficient against their
source rows, verifies the complete row partition and acyclic restoration,
and checks the final sum of squares.  A separate sparse engine expands
the octic completely.  Exact nonzero partial derivatives modulo
$1{,}000{,}000{,}007$ certify that every declared input occurs
in each collected polynomial.  These tests accompany the ordinary proof;
they do not replace its unbounded source calibration.

By Theorem~\ref{thm:acyclic-monic}, each displayed $Q$ has an integer zero
if and only if $P_{248}$ has one.  Compose this elementary, ZFC-formalizable
bijection with the already proved existence equivalence for the fused
front-sign source and its retained C222/ZFC calibration.  For every row
of the table this gives
\[
 \mathrm{ZFC}\ \vdash
 \bigl[\text{$Q$ has no zero in its stated integer domain}\bigr]
 \ \longleftrightarrow\ \operatorname{Con}(\mathrm{ZFC}).
\]
This establishes an equivalence, without assuming or proving either
side.  The new finite algebra is not claimed to be a fresh end-to-end
proof-assistant development.
